To begin, let $\textbf{Choice} = \{C,D\}$, and let $\mathbb{N}$ (the natural numbers) include the source code of all Prisoner's Dilemma algorithms in the sense that each number has a binary representation that may comprise such an algorithm's source code.

Thus, we define the set:
\begin{equation*}
    \textbf{Algorithm} = \{f \mid f \text{ is a } \mathbb{N} \rightarrow \textbf{Choice} \text{ function}\}
\end{equation*}
And a source code function $source: \textbf{Algorithm} \rightarrow \mathbb{N}$.

(Note, we ignore how strategies behave on numbers that do not correspond to an algorithm's source code)